% !TEX program = xelatex
% 贝叶斯公式 · 独立精讲讲义（无课件来源，全文自洽推导）
\documentclass[11pt]{ctexart}
\input{D:/aaa_D_usual/DSH/111学习助手/.dsh/skills/lecture-deepdive/assets/theme-violet}
\settags{贝叶斯公式 · 精讲}{概率与统计}
\setlength{\emergencystretch}{2em}
\setlength{\headheight}{22pt}

\begin{document}

\doctitle{贝叶斯公式}{从「证据的条件概率」到「假设的后验概率」}
{独立精讲，无课件来源：全文每个结论都给出完整推导或可复算的数字\ \textperiodcentered\ 2026-10-23}

%=============================================================
\hdr{§0}{读法、记号与三条主线}

这份讲义只解决一个问题：\textbf{已知证据在假设下的概率，如何反过来求假设在证据下的概率}。
它在教科书里只占半页，却最容易在考试和真实判断里出错，原因不在公式难，而在两个概率的方向被混淆。

\subhdr{记号表（全文只用这一套）}

\par\vspace{2pt}
\noindent\begin{minipage}{\linewidth}
\tabtop
\noindent\begin{tabular}{@{}p{3.2cm}p{3.4cm}p{8.8cm}@{}}
\thead{名字} & \thead{记号} & \thead{一句话含义}\\
\hline
先验概率 & $P(H)$ & 看到证据\textbf{之前}，你认为假设成立的可能性。\\
似然 & $P(E\mid H)$ & 假设成立时，出现这个证据的可能性。注意条件是 $H$。\\
证据概率 & $P(E)$ & 不管假设成不成立，这个证据出现的总可能性；它是归一化常数。\\
后验概率 & $P(H\mid E)$ & 看到证据\textbf{之后}，假设成立的可能性。条件换成了 $E$。\\
似然比 & $\mathrm{LR}$ & $P(E\mid H)/P(E\mid H^{c})$，证据对假设的支持强度。\\
几率 & $\mathrm{Odds}(H)$ & $P(H)/\bigl(1-P(H)\bigr)$，用「几比几」代替「占多少」。\\
\end{tabular}
\tabbot
\end{minipage}

\subhdr{三条主线}

\textbf{第一条（§1）}用一张 10000 人的频数表把两个方向摆在一起，让你先在数字上看清它们不是一回事；
\textbf{第二条（§2--§3）}把它写成公式：条件概率 $\to$ 乘法公式 $\to$ 全概率公式 $\to$ 贝叶斯公式，
并给出几率形式与对数几率形式；
\textbf{第三条（§4）}把事件型的 $H$ 换成连续参数 $\theta$，得到密度形式的贝叶斯更新，并完整演算贝塔--二项共轭。

读完的标准不是记住公式，而是能对任何一个应用题回答：谁是 $H$，谁是 $E$，先验从哪来，分母怎么算。

%=============================================================
\hdr{§1}{为什么「条件要反过来」不是文字游戏}

\subhdr{一张 10000 人的频数表}

设某种病的患病率是 $1\%$，检测的灵敏度（真有病时测出阳性）是 $99\%$，特异度（真没病时测出阴性）是 $95\%$。
换句话说，健康人有 $5\%$ 会被误判为阳性。现在随机抽一个人，测出阳性，他真有病的概率是多少？

把 10000 人按「有病/没病」与「阳性/阴性」两两切开直接数格子：

\par\vspace{2pt}
\noindent\begin{minipage}{\linewidth}
\tabtop
\noindent\begin{tabular}{@{}p{3.4cm}p{3.6cm}p{3.6cm}p{3.4cm}@{}}
\thead{} & \thead{检测阳性 $E$} & \thead{检测阴性 $E^{c}$} & \thead{合计}\\
\hline
真有病 $H$ & $100\times 0.99=99$ & $100\times 0.01=1$ & 100\\
真没病 $H^{c}$ & $9900\times 0.05=495$ & $9900\times 0.95=9405$ & 9900\\
\hline
合计 & $99+495=594$ & $1+9405=9406$ & 10000\\
\end{tabular}
\tabbot
\end{minipage}

阳性的一共 594 人，其中真有病的只有 99 人，于是

\[
P(H\mid E)=\frac{99}{594}=\frac{1}{6}\approx 16.7\% .
\]

而灵敏度说的是 $P(E\mid H)=99/100=99\%$。两个数一个 $99\%$、一个 $16.7\%$，差了近六倍。

\subhdr{差在哪：分子一样，分母不一样}

把两个式子的分数并排写出来就一目了然：

\[
P(E\mid H)=\frac{99}{100},\qquad P(H\mid E)=\frac{99}{594}.
\]

\textbf{分子是同一个 99}——都是「真有病且测出阳性」的那 99 个人；
\textbf{分母换掉了}：前者除以「真有病的 100 人」，后者除以「测出阳性的 594 人」。
所以「条件反过来」在数学上不是把符号对调，而是\textbf{把分母从一个条件组换到另一个条件组}。
分母里多出来的 495 个假阳性，正是把 $99\%$ 拉到 $16.7\%$ 的原因。

\callout{\textbf{自然频数法}：遇到条件概率的判断题，先把人群设成一个大整数（1000 或 10000），
按两条属性切开数格子，再让格子数直接当概率。这个方法不需要记公式，也不会把方向搞反，
它的有效性在认知心理学里有专门的实验证据（Gigerenzer 等人的「自然频数」研究）：
人对计数的直觉远好于对条件概率的直觉。本讲义 §5 例 1 会用公式再算一遍，与这张表互相校验。}

\warnbox{\textbf{别把它当成矛盾}：「灵敏度 $99\%$、后验只有 $16.7\%$」并不违反直觉，
因为阳性这个证据本身很「脏」：594 个阳性里 495 个是健康人被误判。
后验概率的高低由\textbf{先验（患病率）}和\textbf{误判率}共同决定，单看灵敏度什么都说明不了。
这就是所谓「基础率谬误」（base rate fallacy）。}

%=============================================================
\hdr{§2}{从条件概率到全概率公式}

\subhdr{条件概率与乘法公式}

设 $P(B)>0$，在 $B$ 已发生的条件下 $A$ 发生的条件概率定义为

\[
P(A\mid B)=\frac{P(A\cap B)}{P(B)} .
\]

把分母乘过去，就得到\textbf{乘法公式}

\[
P(A\cap B)=P(B)\,P(A\mid B)=P(A)\,P(B\mid A).
\]

最后那个等号只是把 $A$、$B$ 的角色对调再用一次定义，但它极其重要：
\textbf{同一个交事件的概率有两个算法}，一个用 $P(B)$ 打头，一个用 $P(A)$ 打头。
贝叶斯公式就是让这两个算法相等。

\subhdr{完备事件组与全概率公式}

设 $H_{1},\dots,H_{n}$ 两两互斥，且并为全集（也就是「必有一个发生、且只有一个发生」），
则对任意事件 $E$ 有

\[
P(E)=\sum_{i=1}^{n}P(E\mid H_{i})\,P(H_{i}) .
\]

推导只有两步：先把 $E$ 按 $H_{i}$ 切开，$E=\bigcup_{i}(E\cap H_{i})$，且这些交事件两两互斥，
所以概率可加；再对每一项用乘法公式：

\[
P(E)=\sum_{i=1}^{n}P(E\cap H_{i})=\sum_{i=1}^{n}P(E\mid H_{i})\,P(H_{i}) .
\]

\aside{这里用到「$H_{i}$ 两两互斥且并为全集」。如果漏掉某个 $H_{i}$，那一部分人根本没被计入，
$P(E)$ 会算小，后验会算大。应用题里这一步的检查方式是：把所有 $H_{i}$ 的概率加起来，必须等于 1。}

\subhdr{数字例子：先正着算一遍}

某厂有两条产线：A 线产 $60\%$、次品率 $1\%$；B 线产 $40\%$、次品率 $2\%$。
产品混在一起，随机抽一件，它是次品的概率是多少？

这里「哪条线生产」就是完备事件组 $\{A,B\}$，$P(A)=0.6$、$P(B)=0.4$；「次品」记为 $D$，
两个似然分别是 $P(D\mid A)=0.01$、$P(D\mid B)=0.02$。代入全概率公式：

\[
P(D)=0.6\times 0.01+0.4\times 0.02=0.006+0.008=0.014 .
\]

也就是 $1.4\%$。注意它是两条线各自贡献的\textbf{加权和}，权重就是产量占比。
这个 $0.014$ 待会儿会原封不动地出现在贝叶斯公式的分母上。

\subhdr{顺带分清：互斥与独立不是一回事}

初学者最容易把这两个词混用，但它们的含义完全相反。

\textbf{互斥}说的是同一层里两个事件不能同时发生：$P(A\cap B)=0$，于是 $P(A\mid B)=0$。

\textbf{独立}说的是一个事件发生与否不影响另一个的概率：$P(A\mid B)=P(A)$，
等价地 $P(A\cap B)=P(A)P(B)$。

拿一枚硬币：正面与反面互斥，但不独立（出现正面会排除反面）；两次抛掷独立，
但不互斥（完全可以同时出现「第一次正面、第二次正面」）。

数字上更直接：若 $P(A)=0.5$、$P(B)=0.5$ 且互斥，则 $P(A\cap B)=0\ne 0.25=P(A)P(B)$，
所以互斥（在概率都非零时）一定\textbf{不}独立。反过来，独立的两事件一般可以同时发生。

\aside{这个区分在全概率公式里会用到：完备事件组要求的是 $H_{i}$ 之间\textbf{互斥}，
而「两次检测」要连乘似然比，要求的是给定 $H$ 时两次检测\textbf{条件独立}。两处要求不能互换。}

\subhdr{乘法公式的链式形式}

乘法公式可以连用，得到

\[
P(A_{1}\cap A_{2}\cap\cdots\cap A_{n})
=P(A_{1})\,P(A_{2}\mid A_{1})\,P(A_{3}\mid A_{1}\cap A_{2})\cdots P(A_{n}\mid A_{1}\cap\cdots\cap A_{n-1}) .
\]

\textbf{数字例子}：一批产品共 10 件，其中 2 件次品。不放回连续抽 3 件，三件全是正品的概率是

\[
\frac{8}{10}\times\frac{7}{9}\times\frac{6}{8}=\frac{7}{15}\approx 0.4667 .
\]

若按「独立」误算成 $(0.8)^{3}=0.512$，误差来自不放回：第二件的条件概率已经不再是 $0.8$。
这条链式公式也是贝叶斯更新的基础——每来一条新证据，就在条件里多加一项。

%=============================================================
\hdr{§3}{贝叶斯公式：把两个乘法公式摆成等式}

\subhdr{推导只有一行}

对 $H$ 与 $E$，乘法公式给出

\[
P(H\cap E)=P(H)\,P(E\mid H)=P(E)\,P(H\mid E).
\]

只要 $P(E)>0$，把第二个等号两边同除以 $P(E)$：

\[
P(H\mid E)=\frac{P(E\mid H)\,P(H)}{P(E)} .
\]

再把分母用全概率公式对完备事件组 $\{H,H^{c}\}$ 展开，就得到最常用的两事件形式：

\[
P(H\mid E)=\frac{P(E\mid H)\,P(H)}{P(E\mid H)\,P(H)+P(E\mid H^{c})\,P(H^{c})} .
\]

\aside{分母必须是 $P(E)$，也就是「所有可能让这个证据出现的情形」的总和。
分子只算了 $H$ 这一条路，分母要把 $H^{c}$ 那条路也加上——这是贝叶斯公式唯一需要动脑的地方。}

\subhdr{三个因子各干什么}

\par\vspace{2pt}
\noindent\begin{minipage}{\linewidth}
\tabtop
\noindent\begin{tabular}{@{}p{3.0cm}p{4.6cm}p{7.8cm}@{}}
\thead{因子} & \thead{作用} & \thead{它决定什么}\\
\hline
$P(H)$ & 先验：动笔前你的立场 & 若先验很小，再强的证据也拉不动后验（§1 的 $1\%$ 患病率）。\\
$P(E\mid H)$ & 似然：假设成立时证据多容易出现 & 它只描述 $H$ 内部，不含任何关于 $P(H)$ 的信息。\\
$P(E)$ & 归一化：把分子拉回 $[0,1]$ & 它的出现使后验必然合法；漏算它就是把似然当后验。\\
\end{tabular}
\tabbot
\end{minipage}

\subhdr{把 §1 的检测例子用公式算一遍}

$P(H)=0.01$，$P(H^{c})=0.99$，$P(E\mid H)=0.99$，$P(E\mid H^{c})=0.05$。代入：

\[
P(H\mid E)=\frac{0.99\times 0.01}{0.99\times 0.01+0.05\times 0.99}
=\frac{0.0099}{0.0099+0.0495}
=\frac{0.0099}{0.0594}=\frac{1}{6}.
\]

结果与 §1 数格子得到的 $99/594$ 完全一致——因为分子分母同乘了 10000。
这是检查自己有没有算错的好办法：\textbf{频数表与公式必须给出同一个数}。

\subhdr{$n$ 个假设的一般形式}

实际题目里假设往往不止两个。设 $\{H_{1},\dots,H_{n}\}$ 是完备事件组，则对每个 $i$

\[
P(H_{i}\mid E)=\frac{P(E\mid H_{i})\,P(H_{i})}{\displaystyle\sum_{j=1}^{n}P(E\mid H_{j})\,P(H_{j})} .
\]

分母对所有 $j$ 求和，与 $i$ 无关，所以它相当于把每个「贡献」$P(E\mid H_{j})P(H_{j})$ 按比例重新分配。
由此可以得到一条极强的自检：

\[
\sum_{i=1}^{n}P(H_{i}\mid E)=1 .
\]

\textbf{如果算出来的后验加起来不等于 1，一定是分母漏了某一路。}

\subhdr{例题：三条产线的次品来源}

A 线产 $50\%$、次品率 $1\%$；B 线产 $30\%$、次品率 $2\%$；C 线产 $20\%$、次品率 $4\%$。抽到一件次品，它来自哪条线？

\textbf{第一步，正着算分母}（全概率）：

\[
P(D)=0.5\times0.01+0.3\times0.02+0.2\times0.04=0.005+0.006+0.008=0.019 .
\]

\textbf{第二步，反着分配}（贝叶斯）：三路的贡献分别是 $0.005$、$0.006$、$0.008$，
分母 $0.019$ 恰好是它们的和，所以后验就是这三个数各占多少份：

\par\vspace{2pt}
\noindent\begin{minipage}{\linewidth}
\tabtop
\noindent\begin{tabular}{@{}p{3.2cm}p{3.4cm}p{3.4cm}p{3.6cm}@{}}
\thead{产线} & \thead{贡献 $P(D\mid H_{i})P(H_{i})$} & \thead{后验 $P(H_{i}\mid D)$} & \thead{读法}\\
\hline
A & $0.005$ & $5/19\approx 26.3\%$ & 产量最大，但次品率最低，后验反而最小。\\
B & $0.006$ & $6/19\approx 31.6\%$ & 中间。\\
C & $0.008$ & $8/19\approx 42.1\%$ & 产量最小，但次品率是 A 的 4 倍，后验最大。\\
\hline
合计 & $0.019$ & $19/19=1$ & 相加为 1，自检通过。\\
\end{tabular}
\tabbot
\end{minipage}

\textbf{这道题的教学价值}：先验（产量占比）与似然（次品率）在这里给的方向完全相反。
贝叶斯公式不偏向任何一方，它把两者\textbf{相乘}——先验大的被放大一次，似然大的被放大一次，
最后胜出的是乘积大的那个。C 线产量只有 A 的 $40\%$，但次品率是 A 的 4 倍，
乘积一比，C 反而最可疑。

\subhdr{几率形式：把除法变成乘法}

很多场合直接比较「几率」比比较概率方便。定义

\[
\mathrm{Odds}(H)=\frac{P(H)}{1-P(H)} .
\]

对两个互斥完备的假设，后验几率有一个极简洁的形式：

\[
\underbrace{\frac{P(H\mid E)}{P(H^{c}\mid E)}}_{\text{后验几率}}
=
\underbrace{\frac{P(H)}{P(H^{c})}}_{\text{先验几率}}
\times
\underbrace{\frac{P(E\mid H)}{P(E\mid H^{c})}}_{\text{似然比 }\mathrm{LR}} .
\]

推导：把贝叶斯公式分别写 $H$ 与 $H^{c}$ 两遍，分母同为 $P(E)$，相除即消掉。

\textbf{这条式子的意义}：新证据不改变「证据本身的强度」，它只是把先验几率\textbf{乘以}一个固定的倍数。
所以若有两个独立证据 $E_{1},E_{2}$，总似然比就是两个似然比之积，几率可以连乘。

\subhdr{对数几率：连乘变成连加}

对几率取自然对数，乘法就变成加法：

\[
\log\frac{P(H\mid E)}{P(H^{c}\mid E)}
=
\log\frac{P(H)}{P(H^{c})}
+
\log\frac{P(E\mid H)}{P(E\mid H^{c})} .
\]

用 §1 的数字验算一遍（疾病检测）：

\par\vspace{2pt}
\noindent\begin{minipage}{\linewidth}
\tabtop
\noindent\begin{tabular}{@{}p{4.4cm}p{3.4cm}p{7.6cm}@{}}
\thead{量} & \thead{数值} & \thead{说明}\\
\hline
先验几率 & $1/99\approx 0.0101$ & 患病率 $1\%$，即「$1$ 比 $99$」。\\
似然比 $\mathrm{LR}$ & $0.99/0.05=19.8$ & 阳性在病人身上比在健康人身上常见 $19.8$ 倍。\\
后验几率 & $0.0101\times 19.8=0.2$ & 连乘的结果。\\
后验概率 & $0.2/1.2=1/6\approx 16.7\%$ & 由几率换回概率。\\
\hline
$\log$ 先验几率 & $-4.595$ & \\
$\log\mathrm{LR}$ & $+2.986$ & 正号表示这条证据支持 $H$。\\
$\log$ 后验几率 & $-1.609$ & 两项相加；仍为负，说明还不够翻盘。\\
\end{tabular}
\tabbot
\end{minipage}

\warnbox{\textbf{连乘的条件}：$P(E_{1},E_{2}\mid H)=P(E_{1}\mid H)P(E_{2}\mid H)$ 才成立，
即两条证据在给定 $H$ 下条件独立。如果两条证据其实来自同一件事（例如两次都做了同一项检查），
连乘会把同一个信息算两遍，似然比被凭空放大。}

%=============================================================
\hdr{§4}{连续形式：参数也是随机变量}

\subhdr{把「假设」换成「参数」}

事件型的贝叶斯里，$H$ 是离散的几个选项。若未知量是一个连续参数 $\theta$（比如某产品的合格率、
某硬币的正面向上的概率），就要用密度来写：

\par\vspace{2pt}
\noindent\begin{minipage}{\linewidth}
\tabtop
\noindent\begin{tabular}{@{}p{3.4cm}p{4.2cm}p{8.0cm}@{}}
\thead{对象} & \thead{记号} & \thead{含义}\\
\hline
先验密度 & $\pi(\theta)$ & 看到数据前，对 $\theta$ 的信念分布。\\
似然 & $f(x\mid\theta)$ & 参数取 $\theta$ 时，观测到 $x$ 的密度。\\
边缘密度 & $m(x)=\displaystyle\int f(x\mid\theta)\pi(\theta)\dd\theta$ & 事件型里 $P(E)$ 的连续版本。\\
后验密度 & $\pi(\theta\mid x)=\dfrac{f(x\mid\theta)\pi(\theta)}{m(x)}$ & 看到 $x$ 之后的信念分布。\\
\end{tabular}
\tabbot
\end{minipage}

实用上最常写的是它的比例形式，因为 $m(x)$ 与 $\theta$ 无关，只是个常数：

\[
\pi(\theta\mid x)\;\propto\; f(x\mid\theta)\,\pi(\theta) .
\]

这句话可以背成：\textbf{后验 $\propto$ 似然 $\times$ 先验}，比例常数由「后验必须积分为 1」唯一确定。

\subhdr{贝塔--二项共轭：完整演算}

设 $\theta\in(0,1)$ 是某事件的发生概率，先验取贝塔分布

\[
\pi(\theta)=\frac{\theta^{a-1}(1-\theta)^{b-1}}{B(a,b)},\qquad
B(a,b)=\int_{0}^{1}t^{a-1}(1-t)^{b-1}\dd t ,
\]

其中 $a,b>0$ 是给定的超参数，$\mathbb{E}\theta=a/(a+b)$。

现在做 $n$ 次独立试验，成功 $s$ 次、失败 $f=n-s$ 次。似然为二项分布：

\[
f(x\mid\theta)=\binom{n}{s}\theta^{s}(1-\theta)^{f} .
\]

按「后验 $\propto$ 似然 $\times$ 先验」，把与 $\theta$ 无关的因子全部丢掉：

\[
\pi(\theta\mid x)\;\propto\;\theta^{s}(1-\theta)^{f}\cdot\theta^{a-1}(1-\theta)^{b-1}
=\theta^{(a+s)-1}(1-\theta)^{(b+f)-1} .
\]

右边正是 $\mathrm{Beta}(a+s,\,b+f)$ 的核。于是

\[
\pi(\theta\mid x)=\mathrm{Beta}(a+s,\;b+f),
\]

归一化常数随之确定为 $1/B(a+s,b+f)$。共轭的意思是：\textbf{先验与后验属于同一个分布族}，
更新只表现为把两个参数各自加上「成功数」与「失败数」。

\subhdr{后验均值：先验与数据的加权平均}

对 $\mathrm{Beta}(A,B)$ 有均值 $A/(A+B)$，故

\[
\mathbb{E}[\theta\mid x]=\frac{a+s}{a+b+n}
=\frac{a+b}{a+b+n}\cdot\frac{a}{a+b}+\frac{n}{a+b+n}\cdot\frac{s}{n} .
\]

也就是\textbf{先验均值与样本频率（极大似然估计）的凸组合}，权重由「先验样本量 $a+b$」与「数据量 $n$」的比例决定。
这就是为什么 $a+b$ 常被称为\textbf{先验的等效样本量}：

\aside{把 $a,b$ 想成\textbf{伪计数}：$\mathrm{Beta}(3,3)$ 相当于事先见过 3 次成功、3 次失败，
共 6 个虚拟观测，均值 $0.5$。$a,b$ 越大，先验越「顽固」，数据越难把它拉动。}

\subhdr{数字演算：先验 $\mathrm{Beta}(3,3)$，观测 7 成功 3 失败}

\par\vspace{2pt}
\noindent\begin{minipage}{\linewidth}
\tabtop
\noindent\begin{tabular}{@{}p{4.4cm}p{4.4cm}p{6.4cm}@{}}
\thead{量} & \thead{数值} & \thead{说明}\\
\hline
先验 & $\mathrm{Beta}(3,3)$ & 均值 $0.5$，标准差 $\sqrt{0.0357}\approx 0.189$。\\
数据 & $n=10,\ s=7,\ f=3$ & 样本频率 $0.7$。\\
后验 & $\mathrm{Beta}(10,6)$ & 参数相加：$3+7=10$，$3+3=6$。\\
后验均值 & $10/16=0.625$ & 介于 $0.5$ 与 $0.7$ 之间，更靠近数据。\\
权重 & 先验 $6/16=0.375$，数据 $10/16=0.625$ & 验算：$0.375\times0.5+0.625\times0.7=0.625$。\\
后验标准差 & $\sqrt{0.0138}\approx 0.117$ & 比先验的 $0.189$ 小：数据收窄了不确定性。\\
\end{tabular}
\tabbot
\end{minipage}

还有一件事件型里看不到的事：\textbf{先验预测}。在看到数据之前，$s=7$ 的概率是

\[
P(s=7)=\binom{10}{7}\frac{B(10,6)}{B(3,3)}\approx 0.1199 ,
\]

它由贝塔--二项分布给出，用来检查先验是否荒唐：如果先验预测说「几乎不可能出现 7 成功」，
而你偏偏观测到了，那就是先验选得不合适。

\subhdr{归一化常数是怎么定下来的}

上面对共轭只用了「后验 $\propto$ 似然 $\times$ 先验」，但那个 $\propto$ 必须补上一个具体的常数才成为密度。
补法只有一条：\textbf{后验在全空间上的积分必须等于 1}。由贝塔函数的定义直接读出

\[
\int_{0}^{1}\theta^{(a+s)-1}(1-\theta)^{(b+f)-1}\dd\theta=B(a+s,\,b+f),
\]

所以

\[
\pi(\theta\mid x)=\frac{\theta^{a+s-1}(1-\theta)^{b+f-1}}{B(a+s,\,b+f)} .
\]

这也解释了为什么共轭情形这么好用：归一化常数不是新算出来的，而是把参数换成 $(a+s,b+f)$ 之后
\textbf{同一个贝塔函数}。事件型里 $P(E)$ 要老老实实用全概率公式拼；连续共轭里它自动落地。

\subhdr{可信区间：把不确定性写成一个范围}

只要一个数（后验均值）常常不够，还要说清「有多准」。贝叶斯给的是\textbf{可信区间}：
取后验的 $5\%$ 与 $95\%$ 分位数之间的区间，含义是「给定数据后，$\theta$ 落在这个范围内的后验概率是 $90\%$」。

\par\vspace{2pt}
\noindent\begin{minipage}{\linewidth}
\tabtop
\noindent\begin{tabular}{@{}p{3.6cm}p{3.6cm}p{2.6cm}p{5.4cm}@{}}
\thead{后验分布} & \thead{$90\%$ 可信区间} & \thead{宽度} & \thead{说明}\\
\hline
$\mathrm{Beta}(3,3)$（只有先验） & $[0.189,\ 0.811]$ & $0.622$ & 没数据时几乎什么都不知道。\\
$\mathrm{Beta}(10,6)$（$n=10$） & $[0.423,\ 0.809]$ & $0.387$ & 区间收窄约四成。\\
$\mathrm{Beta}(73,33)$（$n=100$） & $[0.613,\ 0.760]$ & $0.147$ & 数据多了十倍，区间再窄一半以上。\\
\end{tabular}
\tabbot
\end{minipage}

\textbf{一个必须说清的差别}：可信区间的概率陈述是「$\theta$ 落在此区间」；
频率派的\textbf{置信区间}不能这么读，它说的是「若重复抽样很多次，约 $90\%$ 的区间会盖住真值」。
两者数值接近但含义不同，考试里常用来区分概念。本讲义所有区间都指前者。

\subhdr{先验怎么选：三条可执行的规则}

\textbf{规则一（有历史数据）}：用历史数据折算出等效样本量，直接定 $a,b$。
例如过去 100 批里 60 批合格、40 批不合格，就取 $\mathrm{Beta}(60,40)$，等效样本量 $100$。

\textbf{规则二（没有任何信息）}：取 $\mathrm{Beta}(1,1)$，也就是 $(0,1)$ 上的均匀分布——
它在任何小区间上的密度都一样，表示「不偏向任何合格率」。
此时后验为 $\mathrm{Beta}(1+7,\,1+3)=\mathrm{Beta}(8,4)$，均值 $8/12=2/3\approx 0.667$，与样本频率 $0.7$ 很接近。

\textbf{规则三（有强信念）}：加大 $a+b$，但要能说清来源。
$a+b$ 越大，先验越难被数据推动；若说不出依据，就不要随便加大。

\aside{还有一个常用的「无信息」选择叫 Jeffreys 先验 $\mathrm{Beta}(1/2,1/2)$，
它在参数变换下形式不变，是理论上更讲究的选择；工程上 $\mathrm{Beta}(1,1)$ 更直观，
两者在数据量大于 20 时几乎无差别。}

%=============================================================
\hdr{§5}{三个例题（统一四步：设事件、写先验、写似然、归一化）}

\subhdr{例 1\quad 疾病检测：与 §1 的频数表互校}

\textbf{题}：患病率 $1\%$，灵敏度 $99\%$，特异度 $95\%$，测出阳性，求真患病概率。

\textbf{解}：设 $H=$ 有病，$E=$ 阳性。先验 $P(H)=0.01$，似然 $P(E\mid H)=0.99$，
健康人误判率 $P(E\mid H^{c})=1-0.95=0.05$。

\[
P(E)=0.99\times0.01+0.05\times0.99=0.0099+0.0495=0.0594,
\]
\[
P(H\mid E)=\frac{0.0099}{0.0594}=\frac{1}{6}\approx 16.7\% .
\]

\textbf{核对}：§1 的频数表给出 $99/(99+495)=99/594$，与上式同值（分子分母同乘 10000 的倒数）。
两个方法独立算、结果一致，说明建模与算术都没错。

\textbf{换个问法}：若再独立做一次同样的检测、仍为阳性，后验变成多少？
此时把刚才的后验 $1/6$ 当作新的先验，再做一次更新：

\[
\frac{0.99\times\frac{1}{6}}{0.99\times\frac{1}{6}+0.05\times\frac{5}{6}}
=\frac{0.165}{0.165+0.0417}\approx 79.8\% .
\]

两次阳性才把概率从 $16.7\%$ 推到约 $80\%$，这就是「一次阳性说明不了太多」的定量版本。

\subhdr{例 2\quad 两工厂次品来源：全概率正着算，贝叶斯反着算}

\textbf{题}：A 线产 $60\%$、次品率 $1\%$；B 线产 $40\%$、次品率 $2\%$。抽到一件次品，它来自 A 线的概率是多少？

\textbf{解}：设 $D=$ 次品。§2 已算过 $P(D)=0.014$。分子是 A 线贡献的那部分 $P(D\mid A)P(A)=0.006$：

\[
P(A\mid D)=\frac{0.6\times0.01}{0.014}=\frac{0.006}{0.014}=\frac{3}{7}\approx 42.9\% ,
\]
\[
P(B\mid D)=\frac{0.4\times0.02}{0.014}=\frac{0.008}{0.014}=\frac{4}{7}\approx 57.1\% .
\]

\textbf{读数}：A 的产量是 B 的 $1.5$ 倍，但次品率只有 B 的一半，所以在「次品」这个条件下，
B 反而更可能。这正是贝叶斯更新在做的修正：\textbf{用似然对先验的分母做再分配}。
两者相加为 1，可作自检。

\subhdr{例 3\quad 连续更新：合格率的贝叶斯估计}

\textbf{题}：某工序的合格率 $\theta$ 未知，先验取 $\mathrm{Beta}(3,3)$；抽检 10 件，7 件合格。求后验分布与后验均值。

\textbf{解}：由 §4 的共轭结论，后验为 $\mathrm{Beta}(3+7,\,3+3)=\mathrm{Beta}(10,6)$，于是

\[
\mathbb{E}[\theta\mid x]=\frac{10}{16}=0.625,\qquad
\mathrm{Var}[\theta\mid x]=\frac{10\times 6}{16^{2}\times 17}\approx 0.0138 .
\]

\textbf{与两种极端做法对比}：只看数据（极大似然）给 $0.7$；只看先验给 $0.5$；
贝叶斯给 $0.625$，是按 $6:10$ 的权重折中。若把抽检量加到 100 件、70 件合格，
后验变成 $\mathrm{Beta}(73,33)$，均值 $73/106\approx 0.689$，已经很靠近 $0.7$——
\textbf{数据越多，先验越不重要}，这就是「先验会被数据洗掉」的定量含义。

\subhdr{例 4\quad 正态--正态：信息按「精度」相加}

\textbf{题}：某物体质量 $\theta$ 未知。先验判断 $\theta\sim N(100,\,10^{2})$；用一台已知测量标准差 $\sigma=15$ 的秤
独立称 9 次，样本均值 $\bar x=110$。求后验均值与后验标准差。

\textbf{解}：正态先验配正态似然仍是正态，结论是\textbf{精度（方差的倒数）相加}：

\[
\frac{1}{\tau_{\text{后}}^{2}}=\frac{1}{\tau_{0}^{2}}+\frac{n}{\sigma^{2}},
\qquad
\mu_{\text{后}}=\tau_{\text{后}}^{2}\left(\frac{\mu_{0}}{\tau_{0}^{2}}+\frac{n\bar x}{\sigma^{2}}\right).
\]

代入 $\mu_{0}=100$、$\tau_{0}^{2}=100$、$\sigma^{2}=225$、$n=9$、$\bar x=110$：

\[
\frac{1}{\tau_{\text{后}}^{2}}=0.01+\frac{9}{225}=0.01+0.04=0.05
\;\Longrightarrow\;
\mu_{\text{后}}=\frac{0.01\times100+0.04\times110}{0.05}=\frac{1+4.4}{0.05}=108 .
\]

后验标准差为 $\sqrt{1/0.05}=\sqrt{20}\approx 4.47$。

\par\vspace{2pt}
\noindent\begin{minipage}{\linewidth}
\tabtop
\noindent\begin{tabular}{@{}p{4.0cm}p{3.4cm}p{7.4cm}@{}}
\thead{量} & \thead{数值} & \thead{含义}\\
\hline
先验权重 & $0.01/0.05=0.2$ & 先验只占两成话语权。\\
数据权重 & $0.04/0.05=0.8$ & 9 次称量顶得上「误差 $15/\sqrt{9}=5$」的一次判断。\\
后验均值 & $108$ & 被数据从 $100$ 拉到 $108$，但没到 $110$。\\
后验标准差 & $4.47$ & 比先验的 $10$ 小了六成，比单次测量的 $15$ 小得多。\\
\end{tabular}
\tabbot
\end{minipage}

\textbf{为什么用「精度」而不是「方差」相加}：精度衡量信息量，两份独立信息合起来，
信息量应当相加，而方差是它的倒数，所以方差反而变小。这与并联电阻的公式同形——
两个电阻并联，总电阻比任何一个都小；两份信息合并，总方差比任何一份都小。

%=============================================================
\hdr{§6}{七个常见错误（每条都配一个能算的反例）}

\subhdr{错误 1：条件倒置}

把 $P(E\mid H)$ 当成 $P(H\mid E)$。§1 的反例：灵敏度 $99\%$，后验只有 $16.7\%$。
两者相等只在极特殊情形下成立（$P(H)=P(E)$ 或 $\mathrm{LR}=1$），一般不等。

\subhdr{错误 2：漏掉基础率}

只写分子 $P(E\mid H)P(H)$，忘了加上 $P(E\mid H^{c})P(H^{c})$。
用 §1 的数字：若把分母写成 $0.0099$，会得到 $P(H\mid E)=1$，显然荒谬。

\subhdr{错误 3：分母用错对象}

分母必须是 $P(E)$，即「所有让 $E$ 出现的情形之和」。若误用 $P(H)$ 作分母，
得到的是 $P(E\mid H)=0.99$，又绕回错误 1。

\subhdr{错误 4：把似然当后验}

「阳性在病人中占 $99\%$，所以阳性者有 $99\%$ 是病人」——这是把 $P(E\mid H)$ 读成 $P(H\mid E)$。
正确的读法是：\textbf{分母换成人，结论才会变}。

\subhdr{错误 5：先验取 0 或 1，永远改不过来}

若 $P(H)=0$，则无论证据多强，

\[
P(H\mid E)=\frac{P(E\mid H)\cdot 0}{P(E)}=0 .
\]

反过来 $P(H)=1$ 时后验恒为 1。这叫 Cromwell 规则：\textbf{概率 0 与 1 不是「很有把握」，而是「拒绝更新」}。
实践中应避免取端点，除非逻辑上真的不可能（例如违反守恒律）。

\subhdr{错误 6：把「后验概率低」读成「证据没有用」}

§1 的检测例子给出 $16.7\%$，很容易被读成「这检测没用」。这是两个不同的问题：

\textbf{证据有没有用}看的是似然比：$\mathrm{LR}=19.8$，意味着阳性者患病的几率是阴性者的 $19.8$ 倍，这很强。

\textbf{后验高不高}看的是先验：患病率只有 $1\%$，所以即使放大近 20 倍，也才到 $16.7\%$。

定量地写：$0.1667/0.01\approx 16.7$，后验是原先的十六倍多。
\textbf{判断证据的价值用 LR，判断该不该担心用后验}——两个问题不要混。

\subhdr{错误 7：跨人群搬用同一个 $P(E)$}

分子里的两个似然 $P(E\mid H)$、$P(E\mid H^{c})$ 是检测本身的属性，可以跨人群搬；
但分母 $P(E)$ 依赖人群的患病率，不能搬。把 $1\%$ 人群算出的 $16.7\%$ 直接安到高危人群（患病率 $20\%$）身上是错的。
用同一组 $0.99$ 与 $0.05$、把先验换成 $0.2$ 重算：

\[
P(H\mid E)=\frac{0.99\times0.2}{0.99\times0.2+0.05\times0.8}=\frac{0.198}{0.238}\approx 83.2\% .
\]

同一个检测、同一个阳性结果，只因人群不同，后验从 $16.7\%$ 变成 $83.2\%$。
\textbf{先验不是「主观偏见」，它是人群信息的载体}；漏掉它，结论一定会偏。

%=============================================================
\hdr{§7}{公式卡与四步流程}

\subhdr{四步流程（任何应用题都按这个顺序写）}

\textbf{第一步 识别事件}：把「假设」与「证据」分别命名，写清哪个是条件、哪个是要问的。
\textbf{第二步 写先验}：$P(H)$ 从题干的比例、人群患病率、历史合格率里读出来。
\textbf{第三步 写似然}：$P(E\mid H)$ 与 $P(E\mid H^{c})$ 分别对应「灵敏度」与「误判率」。
\textbf{第四步 归一化}：分母用全概率公式拼出来，再相除；最后检查各后验之和为 1。

\subhdr{看到这类问法，就用这个式子}

\par\vspace{2pt}
\noindent\begin{minipage}{\linewidth}
\tabtop
\noindent\begin{tabular}{@{}p{6.6cm}p{10.0cm}@{}}
\thead{题面问法} & \thead{直接套的式子}\\
\hline
「测出阳性，真有病的概率？」 & 两事件贝叶斯：分子 $P(E\mid H)P(H)$，分母全概率展开。\\
「次品来自哪条线？」 & 后验按来源分配，各路后验之和为 1。\\
「再做一次检测仍阳性？」 & 把上一次的后验当作这一次的先验，迭代更新。\\
「证据强度怎么比较？」 & 用似然比与几率形式：后验几率 $=$ 先验几率 $\times\mathrm{LR}$。\\
「参数的后验分布？」 & 连续形式：后验 $\propto$ 似然 $\times$ 先验；共轭时只改参数。\\
\end{tabular}
\tabbot
\end{minipage}

%=============================================================
\hdr{§8}{自测 5 题（答案在节末）}

\textbf{题 1}：盒 A 有 3 红 2 白，盒 B 有 1 红 4 白，随机选一个盒子（等可能）再摸一球，得红球。它来自盒 A 的概率？

\textbf{题 2}：A 线产 $70\%$、次品率 $1.5\%$；B 线产 $30\%$、次品率 $3\%$。抽到次品，来自 A 线的概率？

\textbf{题 3}：患病率 $5\%$，灵敏度 $95\%$，特异度 $90\%$。测出阳性，真有病的概率？（用 10000 人频数表核对）

\textbf{题 4}：合格率先验 $\mathrm{Beta}(2,2)$，抽检 10 件 8 件合格。后验分布与后验均值？

\textbf{题 5}：先验 $P(H)=0.1$，似然比 $\mathrm{LR}=5$。用几率形式求后验概率。

\flowerrule

\subhdr{答案与关键一步}

\textbf{答 1}：$\dfrac{0.5\times\frac{3}{5}}{0.5\times\frac{3}{5}+0.5\times\frac{1}{5}}=\dfrac{0.3}{0.4}=0.75$。
关键一步：分母要把「从 B 盒摸到红球」也算进去。

\textbf{答 2}：$P(D)=0.7\times0.015+0.3\times0.03=0.0195$，
$P(A\mid D)=\dfrac{0.0105}{0.0195}\approx 0.538$。
关键一步：A 产量大但次品率低，后验只略高于一半。

\textbf{答 3}：$10000$ 人中患病 $500$，真阳 $475$；健康 $9500$，假阳 $950$；阳性共 $1425$，
故 $P=475/1425=1/3\approx 33.3\%$。关键一步：假阳数目是真阳的两倍，把 $95\%$ 拉到 $1/3$。

\textbf{答 4}：后验 $\mathrm{Beta}(2+8,\,2+2)=\mathrm{Beta}(10,4)$，均值 $10/14\approx 0.714$。
关键一步：先验等效样本量 $4$，数据量 $10$，故后验偏向数据。

\textbf{答 5}：先验几率 $0.1/0.9=1/9$，乘 $\mathrm{LR}=5$ 得后验几率 $5/9\approx 0.556$，
后验概率 $=\dfrac{5/9}{1+5/9}=\dfrac{5}{14}\approx 0.357$。关键一步：几率乘法后别忘了换回概率。

%=============================================================
\hdr{§9}{一页速记}

\begin{center}
\fcolorbox{rulepurple}{bandviolet}{\parbox{0.94\linewidth}{\vspace{4pt}
\textbf{贝叶斯公式}\\[3pt]
\[
P(H\mid E)=\frac{P(E\mid H)P(H)}{P(E\mid H)P(H)+P(E\mid H^{c})P(H^{c})}
\]\\[2pt]
\textbf{几率形式}\quad 后验几率 $=$ 先验几率 $\times$ 似然比\\[2pt]
\textbf{连续形式}\quad 后验 $\propto$ 似然 $\times$ 先验；共轭时参数相加\\[2pt]
\textbf{贝塔--二项}\quad $\mathrm{Beta}(a,b)+(\text{成功 }s,\text{ 失败 }f)\Rightarrow\mathrm{Beta}(a+s,b+f)$\\[2pt]
\textbf{后验均值}\quad $\dfrac{a+s}{a+b+n}$，是先验均值与样本频率的加权平均\\[4pt]
}}
\end{center}

\subhdr{三个必查点}

一、\textbf{方向}：问的是 $P(H\mid E)$ 还是 $P(E\mid H)$？后者是题干直接给的，前者才是答案。
二、\textbf{分母}：有没有把 $H^{c}$ 那条路算进去？两路之和就是 $P(E)$。
三、\textbf{端点}：先验是否取了 0 或 1？取了就永远改不过来。

\flowerrule

\begin{center}
{\small\color{inkgray}本讲义无课件来源，所有结论均给出推导或可复算数字；\\
数值已用 Python 独立复算，与正文一致。}
\end{center}

%=============================================================
\hdr{附录 A}{三组共轭速查（含一个数字算例）}

共轭的意思是「先验与后验同族，更新只改参数」。下面三组覆盖了绝大多数入门题。

\par\vspace{2pt}
\noindent\begin{minipage}{\linewidth}
\tabtop
\noindent\begin{tabular}{@{}p{2.6cm}p{3.4cm}p{3.5cm}p{3.4cm}p{2.3cm}@{}}
\thead{先验} & \thead{似然} & \thead{后验} & \thead{更新规则} & \thead{后验均值}\\
\hline
$\mathrm{Beta}(a,b)$ & 二项：$s$ 成功 $f$ 失败 & $\mathrm{Beta}(a+s,\ b+f)$ & 成功加 $a$、失败加 $b$ & $\frac{a+s}{a+b+n}$\\
$\mathrm{Gamma}(\alpha,\beta)$ & 泊松：$t$ 内 $n$ 次 & $\mathrm{Gamma}(\alpha+n,\ \beta+t)$ & 次数加 $\alpha$、时长加 $\beta$ & $\frac{\alpha+n}{\beta+t}$\\
$N(\mu_{0},\tau_{0}^{2})$ & 正态：$n$ 个观测 & $N(\mu_{\text{后}},\tau_{\text{后}}^{2})$ & 精度相加 & 见 §5 例 4\\
\end{tabular}
\tabbot
\end{minipage}

\textbf{三个算例各一行}：

其一，$\mathrm{Beta}(3,3)$ 加「$7$ 成功 $3$ 失败」得 $\mathrm{Beta}(10,6)$，均值 $0.625$（§4 已完整演算）。

其二，某设备故障率 $\lambda$ 的先验取 $\mathrm{Gamma}(2,1)$（均值 $2$ 次/小时）；
观察 $2$ 小时内发生 $5$ 次故障，后验为 $\mathrm{Gamma}(2+5,\,1+2)=\mathrm{Gamma}(7,3)$，
均值 $7/3\approx 2.333$。与样本频率 $5/2=2.5$ 相比，后验被先验拉低了一点。

其三，先验 $N(100,10^{2})$ 加「$9$ 次称量、均值 $110$、$\sigma=15$」得后验均值 $108$（§5 例 4）。

\aside{参数化约定必须声明，不同教材写法不同：本讲义用 $\mathrm{Gamma}(\alpha,\beta)$ 的\textbf{率参数}形式，
密度 $\propto\lambda^{\alpha-1}e^{-\beta\lambda}$，均值 $\alpha/\beta$；
若用尺度参数 $\theta=1/\beta$，更新规则要改写成 $\beta$ 的倒数形式。抄公式前先确认约定。}

\subhdr{怎么选哪一组}

数据是「几次成功/几次失败」$\Rightarrow$ 贝塔--二项；
数据是「一段时间内发生几次」$\Rightarrow$ 伽马--泊松；
数据是「一堆测量值求均值」$\Rightarrow$ 正态--正态。
选错族不会出错，只是失去共轭，后验要数值积分才能算——这正是共轭的价值所在。

\end{document}
